\documentclass[12pt]{article}
\usepackage[utf8]{inputenc}
\usepackage[T1]{fontenc}
\usepackage{hyperref}
\title{Walmajarri Language Resources for K-12 Curriculum Development}
\author{Compiled by Claude Sonnet 5.5 (claude-sonnet-5-5) for Andr\'as Kornai}
\date{}
\begin{document}
\maketitle

\section{Introduction}
Walmajarri (also Walmatjari or Wolmadjeri; ISO 639-3 wmt, Glottocode walm1241) is a Pama-Nyungan language of the Kimberley region of Western Australia, centred on Fitzroy Crossing, Looma and Lake Gregory. Estimates range from about 500 speakers to about 2,000 (Thieberger 1996). The first language of all children is Kriol. Juwaliny, a closely related variety, is catalogued separately by AIATSIS.

\section{Primary Educational Resources}
According to an Endangered Languages Project profile that draws on older sources, several schools have run Walmajarri programmes: Kulkarriya Community School at Noonkanbah, Fitzroy Crossing State School, Lake Gregory, Billiluna, Gogo and La Grange (for Juwarliny). Current status should be checked with the schools. A practical spelling system developed by Hudson and Richards is in use. Profile: \url{https://endangeredlanguages.com/lang/2578}

\section{Linguistic and Reference Materials}
\begin{itemize}
\item Richards, E. and Hudson, J. (1990), \emph{Walmajarri-English dictionary, with English finder list}, Darwin: Summer Institute of Linguistics. Wikipedia notes that a dictionary is available online: \url{https://wikipedia.com/wiki/ISO_639:wmt}
\item Hudson, J. (1978), \emph{The core of Walmatjari grammar}, Canberra: AIAS. No URL was seen for this.
\item AIATSIS Austlang record for Walmajarri (A66): \url{https://collection.aiatsis.gov.au/node/1507659}
\item AIATSIS Austlang record for Juwaliny (A88): \url{https://collection.aiatsis.gov.au/austlang/language/A88}
\item State Library of Western Australia document that appeared in the search results: \url{https://purl.slwa.wa.gov.au/slwa_b3303747_1180.pdf}. Its contents were not examined.
\end{itemize}

\section{Religious and Community Resources}
None found in the search results.

\section{Digital and Audio Resources}
No audio, video or software resources were found in the search results beyond the online dictionary mentioned by Wikipedia.

\section{Curriculum Development Context}
Walmajarri has existing school programmes in the Fitzroy Valley, a standardised spelling and a published dictionary. These are the most substantial base among the four languages surveyed. Children now speak Kriol as a first language, so a curriculum would be second-language or revitalization oriented. Eirlys Richards worked for sixteen years in Fitzroy Crossing teaching speakers to read their own language, according to a Magabala Books description of \emph{Out of the Desert: Stories from the Walmajarri Exodus}: \url{https://magabala.com/products/out-of-the-desert-stories-from-the-walmajarri-exodus}

\end{document}